\section{Chapter~\ref{sec:motorlearning}. Motor Learning}

The least-mean-squares rule and the advantage of a separate error wire are theorems; the structure of the error-wire circuit, weakening on coincidence, tuning of the trace to the delay, the aftereffect, and spontaneous recovery of a skill have been measured; that the whole circuit works as supervised learning and builds an internal model is the leading hypothesis; the numbers of the two-process model are a calculation with the book's model.

{\footnotesize
\setlength{\tabcolsep}{3pt}\renewcommand{\arraystretch}{1.15}
\begin{longtable}{>{\raggedright}p{0.5cm}>{\raggedright}p{4.0cm}>{\raggedright}p{5.6cm}>{\raggedright}p{2.3cm}>{\raggedright\arraybackslash}p{1.7cm}}
\toprule
No. & Claim & Experiment and what was measured & Source & Status \\
\midrule\endfirsthead
\toprule
No. & Claim & Experiment and what was measured & Source & Status \\
\midrule\endhead
1 & The rule~\eqref{eq:lms} on average follows the gradient of the squared error & A result of adaptive-filter theory: a correction proportional to the input and the output error is stochastic gradient descent on the mean squared error. & \cite{WidrowHoff1960} & theory \\
2 & With an error wire for each output, the speed does not depend on the number of outputs; with one shared number, it falls with the number of noisy neurons (Proposition~\ref{prop:teachers}) & Learning curves of a linear network: learning by random perturbation of the neurons is slower than the exact gradient by a factor equal to the number of perturbed neurons. & \cite{Werfel2005} & theory \\
3 & The error-wire circuit works as supervised learning (Proposition~\ref{prop:errorwire}) & Theories derived from the structure of the circuit. The experiments of rows 7--10 agree with them; that the whole circuit works exactly this way has not been proved. & \cite{Marr1969, Albus1971} & hypothesis \\
4 & The expander: about $7\cdot10^{10}$ small \nt{GLUT} neurons, more than in all the rest of the brain & Counting nuclei in a suspension of dissolved tissue: the human cerebellum has $69\cdot10^9$ neurons out of $86\cdot10^9$ in the whole brain. Stereology: $101\cdot10^9$ small neurons in the cerebellum of elderly men. & \cite{Azevedo2009, Andersen1992cereb} & measured \\
5 & Raising the dimension of the input makes the desired dependence linear & The function-counting theorem: a weighted sum of $n$ inputs with a threshold realizes an arbitrary labeling of about $2n$ vectors in general position. That the cerebellum uses exactly this is part of the hypothesis of row~3. & \cite{Cover1965} & theory \\
6 & About $3\cdot10^7$ output \nt{GABA} neurons, each with on the order of $10^5$ inputs & Stereology of the human cerebellum: $30.5\cdot10^6$ output neurons. Electron microscopy, rat: about $1.75\cdot10^5$ inputs from the expander per output neuron. The inhibitory transmitter of the output neurons is from a review. & \cite{Andersen1992cereb, NapperHarvey1988, Ito2001} & measured \\
7 & Exactly one error wire per output neuron, about once per second, each time a powerful burst & A review of recordings: each output neuron receives one climbing \nt{GLUT} input, which evokes a complex spike at a rate of about $1$~Hz. & \cite{Ito2001} & measured \\
8 & The probability of the burst depends on the direction of the movement error & Monkey, oculomotor cerebellum, saccade adaptation: the direction of the visual error sets the probability of a complex spike, and its size sets the timing; the burst changes the simple spikes on the next trial and shifts the movement along the cell's preferred error. & \cite{Herzfeld2018} & measured \\
9 & Coincidence of the error burst with an active input weakens that input ($-\eta e_i x_j$) & A review of experiments: joint stimulation of an input from the expander and the error wire causes long-term weakening of that input; stimulation of the input alone does not. & \cite{Ito2001} & measured \\
10 & The trace length is tuned to the error delay: with a $100\ms$ delay, the input that was active $100\ms$ earlier is weakened most & Slices and living mice: in the region responsible for oculomotor learning, weakening is narrowly tuned to an interval of about $120\ms$ between input and error, the time the error signal takes in this task; in another region different cells are tuned to different intervals. & \cite{Suvrathan2016} & measured \\
11 & The circuit is an adaptive filter~\eqref{eq:adaptivefilter} that learns an internal model of the body & A model derived from the structure of the circuit. There is no direct experiment in which the learned filter has been measured as the transfer function of the body. & \cite{Fujita1982} & hypothesis \\
12 & Prediction subtracts the consequences of one's own movements, so one cannot tickle oneself & Humans: one's own touch feels less ticklish than the same touch by someone else; in the scanner the somatosensory cortex responds more weakly to one's own touch, and the cerebellum is less active when the movement touches the skin. The explanation by subtraction of a prediction is a hypothesis. & \cite{Blakemore1998} & hypothesis \\
13 & Under an external force the miss shrinks, and after the force is removed the arm misses to the other side & People reach with a robot handle in a force field: the trajectories return to straight lines; after the field is removed, the aftereffect is a near mirror image of the initial distortions, and it transfers to untrained parts of the workspace. & \cite{ShadmehrMussaIvaldi1994} & measured \\
14 & Two processes learn~\eqref{eq:tworate}; after a reverse perturbation the skill returns by itself & Humans, a force field, then a short reverse field and trials with the error clamped: the correction returns by itself toward the old one; a two-process model reproduces this, a one-process model does not. & \cite{Smith2006} & measured \\
15 & The numbers of Fig.~\ref{fig:tworate}: $0.84$ (slow $0.63$) after $200$ movements; $6$ reverse ones; fast $-0.53$, slow $0.46$; rebound to $0.37$ within $40$ movements & Calculation with \texttt{models/\allowbreak motor\_\allowbreak learning.py}, $A_f=0.92$, $B_f=0.1$, $A_s=0.996$, $B_s=0.02$: $0.840$ ($0.634$ and $0.205$), $6$ movements, $-0.531$ and $0.457$, a peak of $0.371$ at the $40$th movement. & \texttt{models/\allowbreak motor\_\allowbreak learning.py} & model \\
16 & Two processes are needed because the body changes at different rates & Theory: optimal estimation with disturbances of several lifetimes yields several filters with different time constants and explains spontaneous recovery and faster relearning. & \cite{Kording2007} & hypothesis \\
\bottomrule
\end{longtable}}
